\chapter{Integration of Elementary Functions}\label{ch:elementary}

We now integrate an element of $E = k(t)[y]/(F)$, where $t$ is a logarithm or
an exponential over $k$, assuming that we can integrate in $k$ and solve
there the differential equations that arise. This is the induction step of
Risch's theorem, which he proved by Puiseux
expansions~\cite[pp.~203--220]{raab2022}, and of Bronstein's version of
it~\cite[Part~2]{bronstein1990}, which uses the Hermite reduction and the
residue criterion first and expansions only where they cannot be avoided.

The chapter is not on the same footing as those of Part~II. For a simple
radical extension the method is complete and rational, and we give it with
proofs and two worked examples. For a general algebraic extension the
sources solve the problem through expansions at the places at infinity and
over $t = 0$; the tutorial says of this, twice, that ``the details are quite
technical''~\cite[\S3.8 and \S3.9]{bronstein1998tutorial}. We state what is
done and where, and do not reproduce it.

Notation is that of Chapter~\ref{ch:monomial}: $w$ is an integral basis of
$E$ over $k(t)$, with derivation matrix $(m_{ij})$ and $e_w$
(Definition~\ref{def:suitable}), and $C$ is the field of constants.

\section{The reduction to a simple integrand}\label{sec:elem-reduction}

Given $f \in E$:
\begin{enumerate}
\item \seealg{HermiteReduce}$(f, w)$ gives $g$ and a simple $h$ with $f = Dg +
  h$ (Proposition~\ref{prop:simple}). The integral of $f$ is elementary if
  and only if that of $h$ is.
\item Compute $R$ by~\eqref{eq:mixed-resultant}. If $R$ does not divide
  $\kappa(R)$, the integral is not elementary
  (Theorem~\ref{thm:residue-criterion}).
\item Otherwise the residues $\tau_P(h)$ at the normal places are constants,
  and we must find $v$, the $c_i$ and the $u_i$
  of~\eqref{eq:liouville-mixed}. By Corollary~\ref{cor:tau-liouville}, $v$ has
  no pole at a normal place.
\end{enumerate}
From here the two cases part.

\section{The logarithmic case}\label{sec:elem-log}

Let $t = \log \eta$ with $\eta \in k^*$, so that $Dt = D\eta/\eta$. Every
finite place is normal, so by Corollary~\ref{cor:tau-liouville} the algebraic
part of the integral of a simple $h$ is integral:
\begin{equation}\label{eq:log-v}
  v = \sum_{i=1}^n B_i w_i, \qquad B_i \in k[t] .
\end{equation}

\paragraph{The logarithms on the curve.} At the finite places the orders of
the $u_i$ are given by the residues, as in
Chapter~\ref{ch:logarithmic-part}. At an infinite place $P$ with
ramification index $r$, Bronstein defines a residue
as follows~\cite[Definition~2.10]{bronstein1990}. Let $a \in \kbar$ be the
coefficient of $t^{-1}$ in the expansion of $h$ at $P$ in powers of
$t^{-1/r}$. If
\begin{equation}\label{eq:tau-infinity}
  a = D\beta + c\, \frac{D\eta}{\eta}
  \qquad \text{with } \beta \in \kbar \text{ and } c \text{ a constant,}
\end{equation}
then $\tau_P(h) = -r c$; otherwise $\tau_P(h)$ is undefined. Since $t$ is a
monomial, $c$ is unique when it exists, and finding it is an integration
problem in $\kbar$: the recursion.

\begin{theorem}\label{thm:log-residues}\standing{cited}{\cite[Propositions~2.11, 2.12 and~2.13]{bronstein1990}}
Let $t$ be a logarithm over $k$ and $h \in E$ simple, with an elementary
integral over $E$. Then $\tau_P(h)$ is defined at every infinite place, and
there is a representation~\eqref{eq:liouville-mixed}, in which the $u_i$ that
lie in $\kbar$ are listed separately,
\[
  h = Dv + \sum_j d_j\, \frac{D\gamma_j}{\gamma_j} + \sum_i c_i\, \frac{Du_i}{u_i},
  \qquad \gamma_j \in \kbar^* ,
\]
such that $\tau_P(h) = \sum_i c_i\, \ord{P}(u_i)$ at every place $P$, finite or
infinite. The sum of the $\tau_P(h)$ over all places, weighted by their
degrees, is zero. If $q_1, \dots, q_l$ is a basis over $\Q$ of the span of the
$\tau_P(h)$ and $\delta_1, \dots, \delta_l$ are the divisors formed from their
coordinates as in~\eqref{eq:delta}, then each $\delta_j$ has a principal
multiple.
\end{theorem}

So the logarithms whose arguments genuinely involve $t$ and $y$ are found as
in Chapter~\ref{ch:logarithmic-part}: residue divisors, now with support at
infinity as well, and a principal multiple of each. The field of constants of
the curve is $\kbar$, which is not a number field, so the test for a
principal multiple needs the specialisation of
Section~\ref{sec:fo-parameters}. Subtracting the logarithms found leaves an
integrand $h_1$ for which we must solve
\begin{equation}\label{eq:E3}
  h_1 = Dv + \sum_j d_j\, \frac{D\gamma_j}{\gamma_j}, \qquad
  v \in E, \quad \gamma_j \in \kbar^*, \quad d_j \text{ constants.}
\end{equation}

\paragraph{The algebraic part.} Substituting~\eqref{eq:log-v}
in~\eqref{eq:E3}, with $h_1 = \sum_i A_i w_i / U$ where $U$ is normal, and
comparing coefficients of the $w_i$, assuming $w_1 = 1$:
\begin{equation}\label{eq:Slog}
  \frac{A_i}{U} = DB_i + \sum_{j=1}^n \frac{m_{ji}}{e_w}\, B_j
     + \epsilon_i \sum_j d_j\, \frac{D\gamma_j}{\gamma_j}
  \qquad (1 \le i \le n),
\end{equation}
where $\epsilon_1 = 1$ and $\epsilon_i = 0$ for $i > 1$. This is the system
$(S_{\log})$ of~\cite[Proposition~2.15]{bronstein1990}: \eqref{eq:E3} has a
solution if and only if~\eqref{eq:Slog} has one with $B_i \in k[t]$. It is a
coupled system of linear differential equations over $k(t)$ for the $B_i$.
In general Bronstein bounds the degrees of the $B_i$
\cite[Proposition~2.14]{bronstein1990} and determines their coefficients by
comparing the expansions of the two sides at the infinite places, following
Risch. When the system is diagonal no expansions are needed.

\begin{theorem}[simple radicals, logarithmic case]\label{thm:log-radical}\standing{proved}{}
Let $t$ be a logarithm over $k$, and let $E = k(t)[z]/(z^n - H)$ be a simple
radical extension with the integral basis $w$
of~\eqref{eq:radical-basis} and the $\lambda_i$
of~\eqref{eq:radical-derivative}, formed with $D$ in place of $d/dx$. Let $h_1
= \sum_{i=0}^{n-1} (A_{i+1}/U)\, w_{i+1}$ with $U$ normal. Then~\eqref{eq:E3}
has a solution if and only if
\begin{enumerate}[label=\textup{(\roman*)}]
\item $A_1/U \in k(t)$ has an integral of the form $B_1 + \sum_j d_j \log
  \gamma_j$ with $B_1 \in k[t]$ and $\gamma_j \in \kbar^*$, and
\item for $1 \le i < n$, the Risch differential equation
  \[
    DB_{i+1} + \lambda_i\, B_{i+1} = \frac{A_{i+1}}{U}
  \]
  has a solution $B_{i+1} \in k[t]$.
\end{enumerate}
The solution is then $v = \sum_i B_i w_i$.
\end{theorem}

\begin{proof}
Since $Dw_{i+1} = \lambda_i w_{i+1}$ and $\lambda_0 = 0$, the first equation
of~\eqref{eq:Slog} reads $A_1/U = DB_1 + \sum_j d_j\, D\gamma_j/\gamma_j$,
and the others read $A_{i+1}/U = DB_{i+1} + \lambda_i B_{i+1}$. And $v$ is of
the form~\eqref{eq:log-v} by Corollary~\ref{cor:tau-liouville}.
\end{proof}

Both are problems over $k(t)$, solved in Volume~I: (i) by the integration
algorithm of \VolI{Chapter~5}, and (ii) by the algorithm for the Risch
differential equation of \VolI{Chapter~6}. This is equations~(23) and~(24)
of~\cite[\S3.8]{bronstein1998tutorial}.

\begin{example}\label{ex:elem-log}
This is Example~14 of~\cite{bronstein1998tutorial}. Consider
\[
  \int \frac{(x^2 + 2x + 1)\sqrt{x + \log x} + (3x + 1)\log x + 3x^2 + x}
            {(x \log x + x^2)\sqrt{x + \log x} + x^2 \log x + x^3}\, dx .
\]
Let $k = \Q(x)$, $t = \log x$ and $E = k(t)[y]/(y^2 - x - t)$. Rationalising
the denominator, the integrand is
\[
  f = \frac{((3x + 1)\,t - x^3 + x^2)\,y - (2x^2 - x - 1)\,t - 2x^3 + x^2 + x}
           {x\,(t + x)(t + x - x^2)} .
\]
\begin{enumerate}
\item With respect to the integral basis $(1, y)$ the denominator $d = x\,(t +
  x)(t + x - x^2)$ is squarefree, so $f$ is simple and the Hermite reduction
  does nothing.
\item $Dd = t^2 + (2 + 4x - 3x^2)\,t + 2x + 2x^2 - 4x^3$, and
  \[
    R(Z) = x^{12} (x - 1)^2 (x + 1)^2 (2x + 1)^2\; Z^3 (Z - 2) .
  \]
  The roots $0$ and $2$ are constants: the criterion is passed.
\item The root $2$ is the residue at the place $P$ where $(t, y) = (x^2 - x,
  -x)$. There is one place $Q$ at infinity, with $r = 2$, and the coefficient
  of $t^{-1}$ in the expansion of $f$ there is
  \[
    a = 1 - 2x + \frac1x = D(x - x^2) + 1 \cdot \frac{Dx}{x} ,
  \]
  so $c = 1$ in~\eqref{eq:tau-infinity} and $\tau_Q(f) = -2$. The residue
  divisor is $\delta = 2P - 2Q$, and it is principal: $\delta = ((x +
  y)^2)$, since $(x + y)(x - y) = x^2 - x - t$ vanishes simply at the two
  points with $t = x^2 - x$ and $x + y$ vanishes at the one with $y = -x$.
\item Subtracting, $h_1 = f - 2\, D(x + y)/(x + y) = (x + 1)\,y / (x\,(t +
  x))$.
\item Theorem~\ref{thm:log-radical}: $A_1 = 0$ gives $B_1 = 0$ and no further
  logarithm; and $\lambda_1 = \frac12 D(x + t)/(x + t) = (x + 1)/(2x\,(t +
  x))$, so the Risch equation is
  \[
    DB_2 + \frac{x + 1}{2x\,(t + x)}\, B_2 = \frac{x + 1}{x\,(t + x)} ,
  \]
  with solution $B_2 = 2$.
\end{enumerate}
Therefore
\[
  \int f\,dx = 2\sqrt{x + \log x} + 2 \log\bigl( x + \sqrt{x + \log x} \bigr) .
\]
The integrand is printed with $x^2 + x + 1$ in place of $x^2 + 2x + 1$
in~\cite[\S2]{bronstein1990}; differentiating the answer shows that the
tutorial's version is the correct one.
\end{example}

\section{The exponential case}\label{sec:elem-exp}

Let $t = \exp \eta$ with $\eta \in k$, so that $Dt = t\, D\eta$. Now the
places over $t = 0$ are special, as are those at infinity, and the algebraic
part can have poles at both. After the Hermite reduction the denominator of
$h$ is $t^m U$ with $U$ normal and $m \ge 0$.

\paragraph{Poles over $t = 0$.} Bronstein removes them one order at a time,
by a reduction that is to the special places what Hermite's is to the normal
ones, except that it can fail. Write $h = \sum_i (b_i^{(0)} w_i)/t^m + \dots$,
where $b_i \in k$ is the value at $t = 0$ of the coefficient of $w_i$ in $t^m
h$; one looks for $c_1, \dots, c_n \in k$ such that $h - D(\sum_i c_i w_i /
t^m)$ has denominator $t^{m-1} U$. Differentiating and evaluating at $t = 0$
gives the system
\begin{equation}\label{eq:Sexp}
  b_i = Dc_i - m\, (D\eta)\, c_i + \frac{1}{e_w(0)} \sum_{j=1}^n m_{ji}(0)\, c_j
  \qquad (1 \le i \le n)
\end{equation}
over $k$, the system $(S_{\exp})$ of~\cite[\S2]{bronstein1990}; here $e_w(0)
\ne 0$ because $e_w$ is normal and $t$ is special.

\begin{theorem}\label{thm:exp-reduction}\standing{cited}{\cite[Proposition~2.18]{bronstein1990}}
Let $t$ be an exponential over $k$ and $h \in E$ simple with denominator $t^m
U$, $m \ge 1$. If $h$ has an elementary integral over $E$,
then~\eqref{eq:Sexp} has a solution in $k$.
\end{theorem}

So either the system has no solution and the integral is not elementary, or
$m$ is lowered. \eqref{eq:Sexp} is a coupled system of $n$ first-order linear
differential equations over $k$. For a simple radical extension it is
diagonal:
\[
  b_{i+1} = Dc_{i+1} + \bigl( \lambda_i(0) - m\, D\eta \bigr)\, c_{i+1}
  \qquad (0 \le i < n),
\]
$n$ separate Risch differential equations over $k$, which is
equation~(26) of~\cite[\S3.9]{bronstein1998tutorial}.

\paragraph{Poles at infinity.} The substitution $\bar t = 1/t$ is again an
exponential over $k$, with $D\bar t = -\bar t\, D\eta$, and it exchanges the
places over $t = 0$ with those at infinity. Running the same reduction for
$\bar t$ removes the poles at infinity. This is the one change of variable
that the mixed case allows.

\paragraph{The rest.} After both reductions, Bronstein shows that $v$ can be
taken in $\kbar$, and the logarithms are found from residues, which now have
to be defined at the special places as
well~\cite[Definition~2.21, Propositions~2.22 to~2.24]{bronstein1990}. That
definition carries an assumption: that $D\eta$ is not a linear combination
with constant coefficients of logarithmic derivatives of elements of
$\kbar$. We have not reproduced this part.

\begin{example}\label{ex:elem-exp}
This is Example~15 of~\cite{bronstein1998tutorial}. Consider
\[
  \int \frac{3\,(x + e^x)^{1/3} + (2x^2 + 3x)\,e^x + 5x^2}{x\,(x + e^x)^{1/3}}\, dx .
\]
Let $k = \Q(x)$, $t = e^x$ and $E = k(t)[y]/(y^3 - t - x)$. The integrand is
\[
  f = \frac{((2x^2 + 3x)\,t + 5x^2)\,y^2 + 3t + 3x}{x\,(t + x)} .
\]
\begin{enumerate}
\item With respect to $(1, y, y^2)$ the denominator $x\,(t + x)$ is normal,
  and $R(Z) = x^8 (1 - x)^3 Z^3$: the criterion is passed and there are no
  residues at the normal places.
\item The denominator is not divisible by $t$, so there is nothing to do
  over $t = 0$. Put $\bar t = 1/t$ and $z = \bar t\,y$, so that $z^3 = \bar
  t^2 (1 + x \bar t)$. By~\eqref{eq:radical-basis} the integral basis is
  $\bar w = (1, z, z^2/\bar t)$, and
  \[
    f = \frac{3x\,\bar t^2 + 3\,\bar t + (5x^2\,\bar t + 2x^2 + 3x)\,\bar w_3}
             {\bar t\,(x^2\,\bar t + x)} ,
  \]
  whose denominator is divisible by $\bar t$: $m = 1$.
\item Here $D\bar t = -\bar t$, so $D\eta = -1$, and the values at $\bar t =
  0$ of the $\lambda_i$ are $0$, $-\frac23$ and $-\frac13$. The diagonal
  system is
  \[
    0 = Dc_1 + c_1, \qquad 0 = Dc_2 + \tfrac13\, c_2, \qquad
    2x + 3 = Dc_3 + \tfrac23\, c_3 ,
  \]
  with solution $c_1 = c_2 = 0$, $c_3 = 3x$.
\item $h = f - D(3x\,\bar w_3/\bar t) = 3/x$, which is in $k$.
\end{enumerate}
Since $\bar w_3 / \bar t = z^2/\bar t^2 = y^2$,
\[
  \int \frac{3\,(x + e^x)^{1/3} + (2x^2 + 3x)\,e^x + 5x^2}{x\,(x + e^x)^{1/3}}\, dx
  = 3x\,(x + e^x)^{2/3} + 3 \log x .
\]
\end{example}

\section{Other routes to the same cases}\label{sec:elem-others}

\paragraph{Trager's reduction for one radical.} The method of
Section~\ref{sec:radical-trager1979} was designed for this
situation~\cite[Appendix~A]{trager1984}: an integrand $R/y$ with $R \in
F(\theta)$, $\theta$ a logarithm or an exponential over $F$, and $y^n \in
F[\theta]$ depending on $\theta$. The reductions at the finite poles are as
before. The reduction of the polynomial part leads, for a logarithm, to a
recursive integration in $F$ for each coefficient, and, when the leading
coefficient of $y^n$ is not a constant, to a first-order linear differential
equation in $F$; for an exponential, to such an equation at each step, both
at $\theta = 0$ and at infinity. Trager shows that each equation has at most
one solution, and that the reductions succeed whenever the integral is
elementary. It is the same computation as
Theorems~\ref{thm:log-radical} and~\ref{thm:exp-reduction} in the diagonal
case, organised around $1/y$ in place of an integral basis, and it finds
the algebraic part only.

\paragraph{Transcendentals over one curve.} Davenport~\cite{davenport1984}
treats the tower in the other order: $K(x, y, \theta_1, \dots, \theta_n)$
with $y$ algebraic over $K(x)$ and each $\theta_i$ a logarithm or an
exponential, transcendental over what precedes it. He shows that the
integration problem and the Risch differential equation are solvable there.
The curve is then fixed, over a field of constants, and the monomials are
added on top of it, so that the transcendental algorithm applies with
coefficients in the function field of the curve. Of the general case, an
algebraic extension over a tower containing transcendentals, he writes that
it ``me semble \^etre non r\'esolu''. That was in 1984; it is the case this
chapter is about, and Bronstein's paper is the first published solution of
it.

\paragraph{Parallel integration.\later} Blake's second preprint on the
Risch--Norman method~\cite{blake2026b} allows the radical ``any position in
the tower'', which it calls ``the setting of Bronstein's algorithm for mixed
elementary functions''. It is a heuristic method with a structure theorem
behind it, not a decision procedure, and we have read only its abstract.

\section{How much of this is an algorithm}\label{sec:elem-status}

For a simple radical over one logarithm or exponential, the procedures above
are rational and complete as far as the algebraic part goes, and the
logarithmic part is that of Chapter~\ref{ch:logarithmic-part} with a larger
field of constants. For a general algebraic extension, Bronstein's own
assessment is the place to
start~\cite[Conclusions]{bronstein1990}. His algorithm computes Puiseux
expansions in the solution of the Risch equation, at infinity in the
logarithmic case, and Laurent expansions at the special places in the
exponential case, and he calls this ``the major obstacle to an efficient
implementation of a complete algorithm''. He names three things that a
rational algorithm would need:
\begin{enumerate}[label=\textup{(\roman*)}]
\item a formula for the residues $\tau_P$ at the special places;
\item a way to solve $Dv = f$ for $v \in k$, with $f$ simple, without the
  expansions of $f$ at the special places;
\item a rational algorithm for the Risch equation.
\end{enumerate}
Chapter~\ref{ch:risch-equation} takes up the third. We know of no source
that supplies the first two in general.

\section*{Exercises}
\addcontentsline{toc}{section}{Exercises}

\begin{exercise}
Use Theorem~\ref{thm:log-radical} to integrate $1/(x \log(x)^{3/2})$
(Example~\ref{ex:mixed-hermite}) starting from the simple integrand that
the Hermite reduction leaves, and $\sqrt{\log x}/x$.
\end{exercise}

\begin{exercise}
Show that $\int \sqrt{\log x}\,dx$ is not elementary: with $y^2 = t$, the
Risch equation of Theorem~\ref{thm:log-radical}(ii) is $DB_2 + B_2/(2xt) =
1$, and it has no solution in $\Q(x)[t]$.
\end{exercise}

\begin{exercise}
In Example~\ref{ex:elem-exp}, verify the expression of $f$ in the basis
$\bar w$ and the three values $\lambda_i(0)$.
\end{exercise}
